\tocdecorline{Hoofdstuk 6 -- Cryptografie}
\section{Les 7 -- Cryptografie}

\begin{center}
    \inlinetext{Vanaf hier niet kennen voor examen}[Magenta!33]
\end{center}

\subsection{Wat is cryptografie? -- Culturele Inleiding}

\subsubsection{Taken van moderne cryptografie}

\subsubsection{Symmetrische en asymmetrische cryptografie}

\subsubsection{Berichten van beperkte lengte en blokcijfers}

\subsubsection{Klassieke substitutiecijfers}

\subsubsection{Poly-alfabetsische systemen en rotormachines}

\begin{center}
    \inlinetext{Vanaf hier terug te kennen voor examen}[Magenta!33]
\end{center}

\subsubsection{DES en AES}

\subsection{Symmetrische cryptografie: AES en het veld \texorpdfstring{$\mathbb{F}_{2^8}$}{F2-8}}

\subsubsection{De vier AES-stappen}

\SimpleHeading{SubBytes}

\SimpleHeading{ShiftRows}

\SimpleHeading{MixColumns}

\SimpleHeading{AddRoundKey}

\subsubsection{Limitaties AES}

\subsection{Publieke-sleutelcryptografie}

\subsubsection{RSA -- publieke sleutel via factorisatie}

\SimpleHeading{Sleutelgeneratie}

\SimpleHeading{Encryptie}

\SimpleHeading{Decryptie}

\subsubsection{Veiligheid van RSA}

\subsubsection{RSA-handtekeningen}

\subsection{Discrete-logaritmeprobleem}

\subsection{Diffie--Hellman-sleuteluitwisseling}

\subsection{ElGamal-encryptie}

\subsection{Schnorr-identificatie}

\subsubsection{Schnorr-identificatieprotocol}

\subsubsection{Betekenis van de eigenschappen}

\subsubsection{Van identificatie naar handtekening}

\subsection{Samenvattende Oefeningen}

\emoji{pushpin} \inlinecode{TODO}[red!22] \emoji{pushpin}

